% 第3章 3.4.1：核 K(x,z)=1+xz（dmu=dz/2）给出的两类方程在各正交方向上的乘子。
% 方向 e_1=1，e_2=sqrt3 x，以及与二者正交的 h_1,h_2,h_3（如 x^2-1/3 等）。
% 第一类 T：乘子 1, 1/3, 0, 0, 0；第二类 I-lambda T（lambda=1/2）：1-1/2=0.5, 1-1/6=0.8333, 1, 1, 1。
\begin{tikzpicture}[font=\figurefont\footnotesize]
\begin{axis}[width=80mm,height=44mm,
  ybar,bar width=3mm,xmin=0.4,xmax=5.6,ymin=0,ymax=1.15,
  axis x line=bottom,axis y line=left,
  xtick={1,2,3,4,5},xticklabels={$e_1$,$e_2$,$h_1$,$h_2$,$h_3$},
  ytick={0,0.3333,0.5,0.8333,1},yticklabels={$0$,$\frac13$,$\frac12$,$\frac56$,$1$},
  ylabel={方向上的乘子},
  ylabel style={font=\figurefont\scriptsize},
  tick label style={font=\figurefont\scriptsize,text=FigureInk},
  ymajorgrids,grid style={draw=FigureGrid,line width=.4pt},
  legend style={at={(0.5,1.03)},anchor=south,legend columns=2,font=\figurefont\scriptsize},
  clip=false]
\addplot[fill=FigureBlue,draw=FigureBlue] coordinates {(1,1) (2,0.3333) (3,0) (4,0) (5,0)};
\addplot[fill=FigureGreenFill,draw=FigureGreen] coordinates {(1,0.5) (2,0.8333) (3,1) (4,1) (5,1)};
\legend{第一类：$T$,第二类：$I-\tfrac12T$}
\draw[FigureMuted,decorate,decoration={brace,mirror,amplitude=3pt},line width=.5pt] (axis cs:2.6,-0.2) -- (axis cs:5.4,-0.2) node[midway,below=3pt,font=\FigureNoteFont\scriptsize,text=FigureMuted] {与$1,x$正交的方向};
\end{axis}
\end{tikzpicture}
